% ================================================================
% ==== FILE: content/m_spaces/mspaces_master.tex
% ================================================================

% ==============================
%  Ontology of Continua — Core
%  M-Spaces: Master Overview
%  FULL MODULE — FINAL
% ==============================

\section{MSpaces}
\label{sec:mspaces-overview}

Die Familie der Metaräume $M_0, M_1, \dots, M_{12}$ bietet die
\emph{Zulässigkeitsstruktur} für Kontinua $K_0\dots K_{12}$.
Jedes $K_i$ kann nur existieren, wenn es vollständig konfiguriert ist
$\Omega(K_i)$ liegt streng im zulässigen Bereich des entsprechenden
Metaraum $M_i$:
\[
\Omega(K_i) \subseteq \Omega(M_i).
\]

M-Räume haben drei universelle Funktionen:

\begin{enumerate}
    \item \textbf{Zulässigkeitsbeschränkungen:}
    Sie definieren, welche Strukturkonfigurationen, Schwellen und Flüsse zulässig sind.
    Sie dienen als globale Regulatoren der Kontinua.

    \item \textbf{Maßangaben:}
    $M_{i+1}$ bestimmt, ob $K_i$ neue Achsen und Übergänge bilden kann
    in ein höherdimensionales Kontinuum.

    \item \textbf{Bauaufsicht:}
    Jeder Metaraum erzwingt die globale Kompatibilität der Operatoren
    wie $\Psi$, $\Phi$, $\Lambda$, $U$ und $\Chi$ und stellt dies sicher
    Kontinua gehorchen den universellen Gesetzen der Evolution, des Zusammenbruchs und der Verzweigung.
\end{enumerate}

M-Räume entwickeln sich nicht „innerhalb“ der K-Hierarchie; stattdessen,
sie bilden eine parallele Aufsichtsstruktur.
Die Beziehung kann wie folgt dargestellt werden:
\[
K_i \hookrightarrow M_i,\qquad
M_i \Rightarrow \text{Einschränkungen für } K_i.
\]

Der Übergang $K_i \to K_{i+1}$ ist nur möglich, wenn:
\[
A(K_{i+1}) \subseteq A(M_{i+1}), \quad
P(K_{i+1}) \subseteq P(M_{i+1}),\quad
\Theta(K_{i+1}) \leq \Theta(M_{i+1}).
\]

Jedes $M_i$ erweitert die Dimensionalität zulässiger Strukturen im Vergleich zu
$M_{i-1}$ stellt die Überstruktur bereit, innerhalb derer sich die nächste Kontinuumsebene befindet
kann entstehen.
